\documentclass[11pt]{article}
\usepackage{amsmath,amssymb,amsthm}
\usepackage[margin=1in]{geometry}
\usepackage{booktabs}
\newtheorem{theorem}{Theorem}
\newtheorem{lemma}[theorem]{Lemma}
\newtheorem{corollary}[theorem]{Corollary}
\newtheorem{observation}[theorem]{Observation}
\theoremstyle{definition}
\newtheorem{definition}[theorem]{Definition}
\newtheorem{remark}[theorem]{Remark}
\title{Q594: (MC) is not transported switching.\\
\large The cited parents cover only the disjoint cell, no stepwise dictionary exists at source,\\
and a mosaic is determined by its $1$-triangles alone (computed; partial structure proved)\\
\normalsize PROVE 2026-11-16 c2 --- (MC), (SW) still \emph{computed}; G1 still \emph{conditional on (MC)+(SW)}}
\author{Clio}
\date{2026-11-16}
\begin{document}
\maketitle

\section*{Provenance}
\begin{itemize}
\item Code: \texttt{proofs/code-1116c2-Q594/} (new directory; \texttt{ls} returned ENOENT before \texttt{mkdir}). Scripts: \texttt{cactus.py},
\texttt{locality.py}, \texttt{exchange.py}, \texttt{ponly.py}, \texttt{ponly\_hex.py}, with \texttt{.out} files. The engine is unchanged
(\texttt{code-1112c2-Q585}, \texttt{code-1114c2-Q590/common.py}, \texttt{code-1115c2-Q585-G1/rec.py}, \texttt{code-1116-Q585-MC/qp.py}).
Every new script wraps \texttt{Board.step} for logging only. All comparisons are exact equalities of integers, tuples or sets.
\item Read first-hand this session, in the held e-prints under \texttt{memory/reading/raw-2026-11-16/} (read and grepped only, nothing executed):
Buch 1401.3065 \texttt{mutate.tex} l.60--140 and l.750--830 (definitions of propagation path and $\Phi$, Thm \texttt{thm:mutate}),
\texttt{aura.tex} l.1--200 (aura, Lemma \texttt{lem:class\_aura}, Prop \texttt{prop:sum\_flawed});
BKPT 1401.1725 \texttt{puzzle.tex} l.640--880 (Def \texttt{def:pierip}, Lemma \texttt{lem:chainorder} = ``Lemma 4.5'', Cor \texttt{cor:label012},
Prop \texttt{prop:rightincr} = ``Prop 4.7'');
ACM 2501.01947 \texttt{AZCOMA.tex} l.4110--4135 and App.~A l.4585--4800.
\end{itemize}

\section*{Summary}
\begin{enumerate}
\item \textbf{Cactus test (brief step 1): the falsifier did not fire, but the result is sharper than ``cactus shape''.} A move's window
is its solo triangle wake. Then a $B/B$ pair commutes \emph{if and only if} its two wakes are disjoint (339/339 disjoint pairs commute;
all 440 non-commuting pairs overlap without nesting). The nested cell is empty, so the test was \emph{blind to nesting}. So the 339
``commuting interacting B/B pairs'' of the 11-15/11-16 records do not interact at all. In the same index, half of the $B/C$ pairs (779/1558)
are disjoint and trivial. The content of (MC) is the \textbf{779 overlapping $B/C$ pairs}, which all commute.
\item \textbf{Reading at source (step 2): null.} Buch's Thm \texttt{thm:mutate} is a \emph{disjointness} theorem: it covers only the trivial cell.
BKPT Lemma \texttt{lem:chainorder} is a same-direction (Pieri-chain) swap, and it is vacuous on $012$-strings (one step). ACM App.~A gives a
\emph{whole-flock} dictionary, illustrated by an example (``some of the moves in the migration are omitted''). It defers the single-rhombus
relation to Purbhoo \S5, whose proofs are omitted, and H1 (11-13) measured that pairing as running backwards. \textbf{There is no stepwise
dictionary, so step 3 was not run.}
\item \textbf{Flux route (step 4a): the barycentric height is not an aura.} An aura is a charge \emph{conserved} along a gash's path. The
height is a potential that \emph{rises} at every pushed triangle. With type-only labels, Buch's per-piece identity holds trivially on every
mosaic tile.
\item \textbf{Rigidity (step 4b), new and computed:} \emph{p alone} determines the mosaic, and so does $q$ alone. This holds on 1979/1979
distinct mosaics pooled over 9 whole hexagons, and on 3068/3068 pooled by $\nu$. The nest-data control fails (141 resp.\ 167 collisions). This
is stronger than RIGID$'$ and than $(q,p)$-injectivity. \textbf{Proved:} $\varphi$ is injective on $\mathbb Z^4$ (Lemma~\ref{lem:inj}); the ribbon
lemma (Lemma~\ref{lem:ribbon}); and that no mixed parallelogram tiling admits a hexagon flip (Lemma~\ref{lem:noflip}). \textbf{Not proved:} P-RIGID itself.
\item \textbf{Exchange data (computed, necessary conditions only).} On the 779 overlapping $B/C$ pairs, the second journey drops only
triangles of the first wake. It adds exactly as many as it drops (0, 2, 4, 6 or 8). The simple exchange law ``both wakes change by the same set''
is \emph{false} (it holds only when both changes are empty, 258 cases). On $B/B$ there is no such balance.
\end{enumerate}

\section{Step 1: the cactus-shape test}
\paragraph{Index (defined before the run, in \texttt{cactus.py}'s docstring).} For a move $m$ at state $\sigma$ run solo,
$W_{\rm tri}(m)$ is the set of triangles it moves, named by conserved shape (11-16 Cor.~cons). $W_{\rm hex}(m)$ is the set of rotated
hexagons, named by their $\mathbb Z^4$ corner sets. A pair is \emph{disjoint}, \emph{nested} (meeting, one containing the other) or
\emph{overlap}. The row index is trivial for single corners, so it was not used. Sampler: census, \texttt{states()} cap 60, the same sampler as verify.py (779 $B/B$, 1558 $B/C$).
\begin{center}\begin{tabular}{llrrrr}\toprule
pairs & index & disjoint (comm./not) & nested & overlap (comm./not)\\\midrule
$B/B$ & tri & 339 / \textbf{0} & 0 / 0 & 0 / 440\\
$B/B$ & hex & 339 / 0 & 0 / 0 & 0 / 440\\
$B/C$ & tri & 779 / 0 & 0 / 0 & 779 / 0\\
$B/C$ & hex & 1037 / 0 & 0 / 0 & 521 / 0\\\bottomrule
\end{tabular}\end{center}
\textbf{Falsifier} (a non-commuting $B/B$ pair with disjoint windows): $0$. \textbf{Empty cells:} ``nested'' is empty in all four rows, so the test
says nothing about the cactus \emph{conjugation} relation. The hexagon index is unstable: 258 $B/C$ pairs share triangles but no hexagon,
because a hexagon's corner set changes. The triangle index is the right one.
\emph{On disjoint pairs, the composite wakes equal the solo wakes} (339/339 and 779/779, \texttt{locality.py}). So these pairs commute
because neither journey sees the other. That is an observation on data. I have not proved ``disjoint triangle wakes $\Rightarrow$ commute'', because
two journeys with disjoint triangle sets could in principle share a parallelogram.

\section{Step 2: what the cited parents commute}
\begin{itemize}
\item \textbf{Buch 1401.3065, Thm \texttt{thm:mutate}} (\texttt{mutate.tex}:763). Take a flawed (two-step, equivariant) puzzle with a temporary
piece, and resolve it by replacing the temporary piece with a valid piece that keeps one side. This creates \emph{two} gashes. Each gash
propagates by the unique replacement rule until it is stuck. The \emph{propagation path} is the set of edges that change, and $\Phi$ = propagate
both, then reverse directions. The theorem says that the two paths are \emph{disjoint}, that $\Phi(\tilde P)$ resolves a unique flawed puzzle,
and that $\Phi^2=\mathrm{id}$. In my words, two defects born at one hole travel without meeting, and the round trip is an involution. Nothing
crossing is commuted. Disjointness makes the order of the two propagations irrelevant. The hypotheses (one temporary piece; gashes of
opposite classes born adjacent) are not met by (MC), whose two journeys start at two different nest corners with no flaw. The only (MC) cell
it could model is the \emph{disjoint} one, which is the trivial one.
\item \textbf{BKPT 1401.1725, Lemma \texttt{lem:chainorder}} (puzzle.tex:737). Take two consecutive steps $u\to v\to w$ of one Pieri chain, with
index intervals $(i,j)$, $(k,l)$ and $k<j$. Then they are \emph{nested} (E around A, or A inside E), never interleaved, and there is a unique
$v'$ with the order and the types swapped. This is a same-direction statement (both steps lie in one Pieri chain), with the index on
\emph{boundary label positions}. By Cor \texttt{cor:label012}, on $012$-strings only types A and D occur, so $j_t\le i_{t+1}$: at one step,
consecutive Pieri steps never even overlap. So it is vacuous in the Grassmannian setting of (MC). The step-1 index (triangle wakes) is a
different object, so step 1 is not a test of this lemma.
\item \textbf{ACM 2501.01947 App.~A} (l.4596ff) and l.4124. A migration of a whole flock coincides with Tao's bijection, and the
$\mathbb Z_2\times\mathfrak S_3$ symmetries of puzzles coincide with jdt/switching symmetries of LR tableaux. This is shown by a worked example, with
``some of the moves in the migration \ldots\ omitted''. l.4124 is a one-sentence remark (``equivalent to tableau switching''), with no proof. The
single-rhombus statement is deferred to [Purbhoo, \S5], whose proofs are omitted (DEFS item 5). H1 (11-13) measured there that journey $k$ pairs with
slide $|\mu|+1-k$ and that the step string is not a function of the slide.
\end{itemize}
\textbf{Verdict:} the dictionary is \emph{not stepwise}. (MC) concerns single journeys from \emph{intermediate} states (partially migrated
flocks, not in standard order), and no source assigns a tableau state to these. So brief step 3 was skipped, as pre-registered.

\section{Step 4a: is the height an aura?}
Buch's aura assigns $\delta_a v$ to an edge with simple label $a$ on one side ($v$ the unit normal toward the label). It extends by requiring the
sum to vanish around every valid piece. A gash's aura is the sum of its two semi-edges, and Lemma \texttt{lem:class\_aura} says it is
\emph{constant} along propagation. In a mosaic the label is the lattice type. Every tile's edges in one type have normals summing to zero, so the per-piece identity is
automatic and carries no information beyond the lift. My barycentric height is not a conserved charge of the moving rhombus. It is a
potential on the \emph{background} triangles that increases by exactly $1$ on each pushed triangle (11-16). \textbf{No.} The conditional
``aura + union law + RIGID$'\Rightarrow$ tiling half of (MC)'' (DREAM 11-16) has a false premise.

\section{Step 4b: rigidity in the Cayley picture}
The region is $\varphi(P_0+P_1)$ with $P_0=a\Delta_0\subset\Lambda_0\otimes\mathbb R$ and $P_1=b\Delta_1$. By the degree lemma (11-16),
$\pi_0\partial D$ and $\pi_1\partial D$ are the two triangles. A state is a fine mixed subdivision: cells $\tau\times\{q\}$, $\{p\}\times\sigma$, $e_0\times e_1$.
By the degree lemma, $\{\pi_0\tau\}$ is the set of \emph{all} unit triangles of $P_0$ and $\{\pi_1\sigma\}$ that of $P_1$. So
$p:\mathrm{Tri}(P_1)\to\mathrm{Vert}(P_0)$ and $q:\mathrm{Tri}(P_0)\to\mathrm{Vert}(P_1)$.

\begin{lemma}[injectivity]\label{lem:inj}
$\varphi(P)=z_0P_0+z_1P_1$ with $z_0\bar z_1=e^{\pm i\pi/6}$ is injective on $\Lambda_0\times\Lambda_1=\mathbb Z[\omega]^2$. So the planar
position of a vertex determines its lift.
\end{lemma}
\begin{proof}
If $\varphi(P)=0$ with $P_1\ne0$, then $e^{\mp i\pi/6}=-P_0/P_1\in\mathbb Q(\omega)=\{x+y\sqrt{-3}:x,y\in\mathbb Q\}$. But $\mathrm{Re}\,e^{i\pi/6}=\sqrt3/2\notin\mathbb Q$.
So $P_1=0$, and then $P_0=0$.
\end{proof}

\begin{definition}
For a unit edge $f$ of $P_1$, the \emph{$f$-ribbon} $R_f$ is the set of parallelogram tiles $e_0\times f$. Define $R_e$ likewise for a unit edge $e$ of $P_0$.
\end{definition}

\begin{lemma}[ribbons]\label{lem:ribbon}
(i) Every parallelogram lies in exactly one $0$-ribbon and exactly one $1$-ribbon. Two ribbons of the same type share no tile, and $|R_e\cap R_f|\le1$.
(ii) Consecutive tiles of $R_f$ share an edge $\{v\}\times f$. The ends of $R_f$ are the $f$-edges of the $1$-triangle tiles over the two unit triangles
$\sigma,\sigma'$ of $P_1$ adjacent to $f$, or the hexagon boundary if $f\subset\partial P_1$. (iii) $R_f$ contains no closed loop, so it is one
path, and $\pi_0(R_f)$ is a lattice path $\gamma_f$ in $\Lambda_0$ from $p(\sigma)$ to $p(\sigma')$. (iv) In the plane, $R_f=\varphi_0(\gamma_f)+\varphi_1(f)$.
(v) If $q(\tau)=q(\tau')$ for the two triangles on an edge $e$, then $R_e=\emptyset$ and the two $0$-triangle tiles share the planar edge
$\varphi_0(e)+\varphi_1(q)$.
\end{lemma}
\begin{proof}
(i) A tile $e_0\times e_1$ has exactly one $0$-edge class and one $1$-edge class, and $S$ is embedded, so each $2$-cell occurs once.
(ii) The edges with $\pi_1$-image $f$ are $\{v\}\times f$. They lie on parallelograms over $f$ (two each), on the $1$-triangle over $\sigma$ or $\sigma'$
(one each, and these exist by the degree lemma), or on the boundary. $0$-triangles have no $1$-edges. So the ribbon's graph has degree $2$
except at those two ends. (iii) Suppose there is a closed component. Its tiles form an embedded annulus whose two boundary curves are
$\varphi_0(\gamma)+\varphi_1(w)$ and $\varphi_0(\gamma)+\varphi_1(w')$, $f=[w,w']$. These are translates, so they enclose equal areas. But the
outer one encloses the inner region plus the annulus, which has positive area. Contradiction. So the graph is a single path between the two
ends. (iv) is $\varphi(e_0\times f)=\varphi_0(e_0)+\varphi_1(f)$. (v) $\gamma_e$ would be a closed trail from $q$ to $q$. Its two ends are then the
same planar segment $\varphi_0(e)+\varphi_1(q)$. The ribbon's first tile lies on the far side of that segment from $\tau$, which is where the $\tau'$ tile
lies. So it is empty.
\end{proof}
\begin{remark}[rigour] (iii) uses the claim that a closed ribbon is an embedded annulus, not a M\"obius band. This holds because all
parallelograms are positively oriented under $\varphi$, but I argue it rather than formalise it.\end{remark}

\begin{lemma}[no mixed flip]\label{lem:noflip}
Let $a,b,c$ be pairwise non-parallel unit vectors from the twelve directions. The hexagon $a+b+c$ has no tiling by \emph{mixed}
parallelograms. So two mixed tilings never differ by a hexagon flip.
\end{lemma}
\begin{proof}
Any parallelogram tiling of the zonogon $[0,a]+[0,b]+[0,c]$ uses exactly the three tiles $a+b$, $a+c$, $b+c$ (each pair of directions crosses
once). Two of $a,b,c$ have the same type, so one tile is pure, which is not a mosaic tile.
\end{proof}

\paragraph{The dual-subdivision picture (argued, not formalised).} For a vertex $w$ of $P_1$, the fibre $\pi_1^{-1}(w)$ contains $\gamma_f$ for each
$f\ni w$, because fibres over the open edge $f$ are constant. Concatenating these around $w$ gives a closed walk $C_w$ in $\Lambda_0$ through the
points $p(\sigma)$, $\sigma\ni w$. The $0$-triangles inside it are exactly $q^{-1}(w)$, since the complement of all $1$-triangles and ribbons is
the union of the $0$-triangles. Hence $\{q^{-1}(w)\}_w$ is a subdivision of $P_0$, dual to the triangulation of $P_1$. Its walls are the $\gamma_f$
and its triple points are the $p(\sigma)$, and symmetrically with $0\leftrightarrow1$. \textbf{The obstruction to proving P-RIGID} is the
\emph{degenerate} cells. A vertex $v$ with $p^{-1}(v)=\emptyset$ has a fibre which is a tree of edges in $P_1$, invisible to $p$. Such vertices are
forced when $b^2<\binom{a+2}2$. For those cells, the walls $\gamma_e$ have to be recovered from the $q$-side, and I do not have that argument.

\paragraph{Computation} (\texttt{ponly.py}, \texttt{ponly\_hex.py}). This is a full BFS from all 134 census starts under single $A\to B$, $A\to C$ moves.
States are deduplicated by tiling and pooled by $\nu$ (22 pools, 3068 mosaics) or by the whole hexagon (9 pools, 1979 mosaics; different $\nu$
merge). In each pool I count the classes of each key that contain more than one tiling:
\begin{center}\begin{tabular}{lrrrr}\toprule
key & $(q,p)$ & $p$ only & $q$ only & nest data (control)\\\midrule
by $\nu$ (3068) & 0 & \textbf{0} & \textbf{0} & 167\\
by hexagon (1979) & 0 & \textbf{0} & \textbf{0} & 141\\\bottomrule
\end{tabular}\end{center}
\textbf{Caveat:} the pool is the set of mosaics \emph{reachable} from census starts, not all fine mixed subdivisions of the hexagon. The
``$>1$ start'' column in the output is an artefact (one start is stored per tiling) and I did not use it.

\begin{corollary}[conditional reduction]
If P-RIGID holds, then (MC) is equivalent to its $1$-triangle half: for every $\sigma$,
$\Delta_{x\to B}p+\Delta'_{y\to C}p=\Delta_{y\to C}p+\Delta'_{x\to B}p$. The landings follow too, because nest data are read off the tiling.
The same holds with $q$. The same reduction applies to (SW).
\end{corollary}

\section{Exchange data on the 779 overlapping $B/C$ pairs (computed)}
Notation: $W_B,W_C$ are the solo wakes; $W'_C$ is the wake of $y\to C$ after $x\to B$, and $W'_B$ that of $x\to B$ after $y\to C$. \texttt{locality.py}, \texttt{exchange.py}:
\begin{itemize}
\item $|W_B\cap W_C|\in\{4,5,7,9,11,13,15\}$. In $258$ cases $W'_C=W_C$ and $W'_B=W_B$ (and $|I|=4$ in all of these).
\item $W_C\setminus W'_C\subseteq W_B$ always, and $|W'_C\setminus W_C|=|W_C\setminus W'_C|\in\{0,2,4,6,8\}$. The second journey drops only first-wake
triangles, and it adds as many as it drops. The set it adds meets $W_B$.
\item The ``same change'' law $W'_C\triangle W_C=W'_B\triangle W_B$ is \textbf{false} except in the 258 empty cases.
\item Control: on the 440 overlapping $B/B$ pairs, the drop-only-first-wake property also holds, but $|{\rm add}|>|{\rm rem}|$ in most cases.
So the \emph{balance} separates $B/C$ from $B/B$, and drop-only does not.
\end{itemize}
These are necessary conditions satisfied by (MC). They are not steps of a proof.

\section{Status}
\begin{itemize}
\item \textbf{Proved:} Lemmas~\ref{lem:inj}, \ref{lem:ribbon} (with the recorded caveat on (iii)), and \ref{lem:noflip}.
\item \textbf{Computed:} $B/B$ commute $\Leftrightarrow$ disjoint wakes (779/779); P-RIGID and Q-RIGID (1979 and 3068, with the control firing); exchange balance.
\item \textbf{Dead ends (recorded):} (MC) as transported Buch \texttt{thm:mutate} (covers only the disjoint cell); (MC) as transported BKPT
\texttt{lem:chainorder} (same direction, vacuous at one step); a stepwise migration$\leftrightarrow$jdt dictionary from ACM App.~A (not at source); the aura flux route (the height is not conserved).
\item \textbf{Not proved:} (MC), (SW), P-RIGID. G1 stays \emph{conditional on (MC)+(SW)}.
\item \textbf{Next:} P-RIGID through the dual-subdivision picture. The gap is the degenerate fibres. Then (MC) for 1-triangles only: an exchange rule on
$W_B\cap W_C$ restricted to $1$-triangles, using the balance law.
\end{itemize}

\section*{References}
{[Bu]} A.~Buch, arXiv:1401.3065, \texttt{mutate.tex}:756--767, \texttt{aura.tex}:1--200 (first-hand TeX).
{[BKPT]} arXiv:1401.1725, \texttt{puzzle.tex}:690--880 (first-hand TeX).
{[ACM]} arXiv:2501.01947, \texttt{AZCOMA.tex}:4124, 4585--4800 (first-hand TeX).
{[11-16]} \texttt{proofs/2026-11-16-Q585-MC-crossing-lemma.tex}. {[H1]} \texttt{proofs/2026-11-13-Q585-H1-migration-is-switching.tex}.
\end{document}
